\section{Artifact, API, and Release Acceptance}
\label{app:acceptance}

This appendix summarizes the normative acceptance boundary for the software
artifact and its public API. It introduces no additional measurement and does
not replace the digest-bound evidence described in
Section~\ref{sec:evaluation}.

\subsection{Artifact boundary}

An accepted release is built from one reviewed, clean \code{main} commit and
identifies the deployed image by its complete immutable GHCR digest. A branch,
local image ID, candidate tag, semantic-version tag, or \code{latest} alias is
not a deployment identity by itself. Aliases may be published only when they
resolve to the already accepted digest without rebuilding the image.

The runtime artifact is limited to the pinned \modelname{} checkpoint
and the decoder-only speech-tokenizer material required for inference. The
image embeds those weights and must run without a host model mount. Voice
cloning, reference-audio conditioning, speaker enrollment, the speech-tokenizer
encoder, Base and CustomVoice checkpoints, and the 0.6B model remain outside
the artifact contract.

\subsection{API acceptance surface}

Table~\ref{tab:appendix-api} lists the complete public endpoint surface used by
release acceptance. The canonical synthesis endpoint accepts text and a
natural-language VoiceDesign description. Progressive output uses
\code{multipart/mixed} with 24~kHz mono signed 16-bit PCM; buffered output uses
a bounded RIFF/WAVE response.

\begin{table}[h]
  \centering
  \caption{Public endpoints covered by the release acceptance boundary.}
  \label{tab:appendix-api}
  \small
  \begin{tabularx}{\textwidth}{>{\raggedright\arraybackslash}p{0.34\textwidth}X}
    \toprule
    Endpoint & Acceptance purpose \\
    \midrule
    \code{GET /health/live} & Process and HTTP-loop liveness. \\
    \code{GET /health/ready} & Shared-engine readiness after the complete native warm-up. \\
    \code{GET /v1/capabilities} & Effective VoiceDesign-only languages, formats, and limits. \\
    \code{POST /v1/voice-design/speech} & Canonical progressive PCM or buffered WAV synthesis. \\
    \code{POST /v1/audio/speech} & Narrow buffered-WAV compatibility behavior. \\
    \code{DELETE /v1/requests/\{request-id\}} & Bounded cancellation of an admitted request. \\
    \code{GET /metrics} & Prompt-free request and engine-health metrics. \\
    \bottomrule
  \end{tabularx}
\end{table}

API acceptance requires readiness, progressive delivery before synthesis
completion, a valid buffered WAV response, cancellation, bounded shutdown and
restart behavior, and payload-free operational metrics. The compatibility
endpoint is not treated as a complete implementation of a third-party audio
API. Public deployment additionally requires an authenticated, rate-limited
TLS-terminating proxy because the native service does not provide those edge
controls.

\subsection{Evidence and reproducibility gates}

Publication requires the complete schema-v1.2 comparison: Native and stock
SGLang at B1, B3, and B6 in two rounds, with the shared workload and client,
checksum-inventoried raw evidence, the declared warm-up minimum, and at least
200 successful measured requests in every cell. Native successes require
explicit natural end-of-sequence. Stock SGLang termination remains unknown and
must pass the documented exclusive length-boundary gate without being relabeled
as natural EOS. Telemetry must retain the documented cadence, measured-phase
boundaries, and zero-competing-CUDA-process requirement.

The production manifest, benchmark report, paper data, figures, and tables must
all derive from that same validated evidence root. Release acceptance then
binds the clean source commit, embedded model identity, immutable OCI digest,
SBOM, provenance, vulnerability and secret scans, keyless signature, anonymous
empty-store pull, and digest-specific GPU checks. The semantic-version and
\code{latest} aliases are promoted only after those gates pass and must continue
to resolve to the accepted digest.

A reproducer should record the source tag and commit, immutable image digest,
model revision, evidence-manifest SHA-256, workload SHA-256, hardware, driver,
and CUDA versions. The manifest must be validated before regenerating the
report or paper data. Any change to hardware, workload, comparator, sampling,
duration, or termination semantics must be disclosed as a different experiment
rather than presented as a reproduction of the release result.
